\documentclass[11pt]{article}
\usepackage[T1]{fontenc}
\usepackage[utf8]{inputenc}
\usepackage{lmodern}
\usepackage[margin=1in]{geometry}
\usepackage{amsmath,amssymb,amsthm,mathtools,microtype,booktabs}
\usepackage[hidelinks]{hyperref}
\newtheorem{theorem}{Theorem}[section]
\newtheorem{lemma}[theorem]{Lemma}
\newtheorem{corollary}[theorem]{Corollary}
\newcommand{\Pow}{\mathcal P_+}
\newcommand{\ow}[1]{\mathbin{\star_{#1}}}
\newcommand{\A}{\mathbf A_I^{\mathrm{sym}}}
\newcommand{\B}{\mathsf B}
\title{Owned Choice with Role Switching:\\Pair Tests and a Partial-Permutation Normal Form}
\author{KFOL research working note}
\date{28 October 2026}
\begin{document}
\maketitle
\begin{abstract}
We retain internal role-switching in the intended signature of rational
owned choice. For any number of roles, every failed term equation has a
countervaluation assigning at most two possible winners to each variable.
Every singleton and pair is necessary, so the unique minimum fixed input
test domain has $K(K+1)/2$ values. For three roles, we classify all unary
term operations: there are exactly $751$. On ambiguous inputs their
top-preserving part acts by partial permutations, with unrestricted
singleton rows subject to explicit bounds and mandatory images. The
resulting $745$-element monoid maps onto the $34$-element symmetric inverse
monoid. Unary synthesis in the no-switching reduct and the necessity
direction of its test-domain theorem are now kernel-checked. The new
role-switching results have mathematical proofs and independent finite
cross-checks, not full Lean proofs yet.
\end{abstract}

The follow-up note \texttt{general.tex} identifies the ambiguity-action
quotient for every $K$, using an affine extension of the known
factorpower of $\operatorname{Sym}(I)$. The partial-permutation monoid
below is exactly its three-role case. That note also establishes
coNP-complete equational validity and the unary units at every $K$.

\section{The intended signature}

Fix a finite role set $I$, $|I|=K\geq3$. Values are nonempty subsets of
$I$. They record possible winners over full pure subgame-perfect
equilibrium profiles of finite public games with a unique winner at
each leaf. A role values its own win at one and all other winners at
zero. The operations are
\[
 S\ow r T=\begin{cases}\{r\},&S=\{r\}\text{ or }T=\{r\},\\
 S\cup T,&\text{otherwise},\end{cases}
 \qquad S\longmapsto\pi S\quad(\pi\in\operatorname{Sym}(I)),
\]
together with the winner constants $\{i\}$. Syntactically,
\[
 t ::= \langle i\rangle\mid x_j\mid t\oplus_r t\mid\pi(t).
\]
An internal role switch renames owners and winners throughout its
component. Its valuation is $\pi(\mathsf{val}(t))$, with variables
interpreted in the original environment. It is retained as part of the
intended logic, as in the original proposal and the known owner-indexed
precedents\cite{Abramsky2007,OldeLoohuisVenema2010}.

The smaller signature without internal switches is a reduct used to
analyze this algebra. Its $K+1$-state one-player quotients and $K+2$-input
test domains do not apply unchanged here. With internal switches, every
player's contextual possible-win observation distinguishes every value:
if $S,T$ differ at $i$, a permutation taking $i$ to the observing role
separates them. This privileges no observing role; each can perform the
same transport.

The earlier notes \texttt{report.tex} and \texttt{structure.tex} supply
the full-strategy interpretation, unary synthesis in the reduct, and the
congruence classifications. In particular, the full-signature
congruences form the chain of ambiguity truncations
$c_m(S)=S$ for $|S|\leq m$, and $c_m(S)=I$ otherwise. These are uniform
matroid closures\cite{Oxley2011}. The full algebra is subdirectly
irreducible. The reduct's product representation is a \emph{subdirect}
product, generally not the whole Cartesian product.

Ordinary union is definable using three distinct owners:
\[
 J(x,y)=(x\ow r y)\ow t(x\ow s y)=x\cup y.
\]
All nonempty constants are therefore ground definable. Throughout, terms
use only the displayed constructors; unions and larger constants below
are these derived terms.

\section{Every equation has a singleton-or-pair countervaluation}

Assign each input variable a status: a specified singleton or an
unspecified ambiguous set. Fixing the statuses fixes every test for
whether an intermediate result is a singleton. The remaining dependence
has a restricted positive form.

\begin{lemma}[Permutation provenance on a status cell]\label{lem:provenance}
On every fixed status cell, a term denotes
\begin{equation}\label{eq:provenance}
 C\ \cup\ \bigcup_j\ \bigcup_{\pi\in P_j}\pi S_j,
\end{equation}
where $C\subseteq I$, only ambiguous input variables occur in the union,
and every $P_j$ is a finite, possibly empty, set of permutations.
\end{lemma}
\begin{proof}
A winner or singleton-status input contributes a constant. An ambiguous
input contributes its identity permutation. A role switch transports
$C$ and composes the permutations. A nonempty $P_j$ contributes an
ambiguous set, so that expression cannot equal an owner singleton. Thus
an owned choice collapses exactly when some child is a constant owner
singleton on the cell. Otherwise union its constants and coefficients.
This proves the representation by induction. It is not asserted to be
canonical: distinct coefficients may have the same action on ambiguous
inputs.
\end{proof}

\begin{theorem}[Pair countervaluations]\label{thm:pairs}
If an equation between terms in this full signature fails, it fails at a
valuation in which every variable has one or two possible winners.
\end{theorem}
\begin{proof}
Take a failed valuation, fix its statuses, and choose a role $k$ present
in the left output and absent from the right, interchanging sides if
necessary. Apply Lemma~\ref{lem:provenance} to both terms.

If $k$ lies in the left constant, it survives every valuation in the
cell. Otherwise it has a witness $k=\pi(i)$ for one ambiguous input
$S_j$, one coefficient permutation, and some $i\in S_j$. Replace that
input by a pair contained in it and containing $i$. Replace every other
ambiguous input by any pair contained in it. Leave singleton inputs
fixed. The left witness remains. The right output can only shrink under
these replacements, since every coefficient is a permutation image and
the status cell is unchanged. It still excludes $k$.
\end{proof}

\begin{lemma}[Resetting atomic inputs]\label{lem:reset}
If a unary term function $F$ is extensive, repeatedly replacing its
current value by $F(x)\ow i x$, for each $i\in I$, gives a term
$\operatorname{Reset}(F)$ which fixes all singletons and agrees with $F$
on every ambiguous input.
\end{lemma}
\begin{proof}
At input $\{i\}$, its own suppression returns $\{i\}$. A different
owner preserves the current value because it contains $i$. On ambiguous
input $S$, the current value contains $S$; neither operand is a
singleton, and their union is still the current value. Extensivity is
preserved at each step.
\end{proof}

\begin{theorem}[Unique minimum fixed test domain]\label{thm:optimal}
A fixed input domain tests every term equation, at every arity, exactly
when it includes every singleton and every two-element set. Consequently
the unique minimum domain is
\[
 D_2=\{S\subseteq I:1\leq|S|\leq2\},\qquad |D_2|=K(K+1)/2.
\]
\end{theorem}
\begin{proof}
Sufficiency is Theorem~\ref{thm:pairs}. Each singleton is necessary
already in the reduct: its unary synthesis gives a term differing from
the constant $I$ only at that singleton. These separators are now
kernel-checked in \texttt{Synthesis.lean}.

For necessity of a particular pair $P$, choose $k\notin P$, put
$C=I\setminus\{k\}$, and form the extensive term
\[
 F(x)=x\cup C\cup\bigcup_{i\in I\setminus(P\cup\{k\})}(i\ k)(x).
\]
Here $(i\ k)$ is a transposition. If $x\subseteq P$, then $F(x)=C$.
Otherwise $x$ contains $k$, or contains an $i$ whose transposition
supplies $k$, and $F(x)=I$. Thus, on ambiguous inputs, $F$ differs from
$I$ only at $P$. By Lemma~\ref{lem:reset},
$\operatorname{Reset}(F)$ and $\operatorname{Reset}(I)$ agree at every
singleton and every ambiguous input other than $P$, and disagree at
$P$. Every complete input domain must therefore contain $P$.
\end{proof}

All $K$ owners and role switches remain available in these tests.
Intermediate outputs may contain any number of winners. The domain
$D_2$ is not closed under the operations; it is an input test domain,
not a replacement algebra. For four roles it has ten values, while the
reduct admits six. The pair separator above explicitly refutes reusing
that six-value test domain after adding internal role switches.

The theorem supplies a decision procedure, uniform in $K$, for term
equality: test $D_2^n$. It does not by itself supply a finite equational
basis or a new proof calculus.

\section{Three roles: the exact unary normal form}

Now let $I=\{a,b,c\}$. A partial permutation $p$ is an injective partial
map from $I$ to itself. Put
\[
 C_p=I\setminus\operatorname{ran}(p),\qquad
 F_p(S)=C_p\cup p[S\cap\operatorname{dom}(p)].
\]

\begin{theorem}[All unary term operations]\label{thm:unary}
Every unary term function is either a constant or has the following
unique normal form:
\begin{equation}\label{eq:unary}
 f(S)=\begin{cases}D_i,&S=\{i\},\\F_p(S),&|S|\geq2,\end{cases}
\end{equation}
where $p$ is a partial permutation and the nonempty rows satisfy
\[
 \begin{array}{ll}
 \{p(i)\}\subseteq D_i\subseteq C_p\cup\{p(i)\},
       &i\in\operatorname{dom}(p),\\
 \varnothing\ne D_i\subseteq C_p,
       &i\notin\operatorname{dom}(p).
 \end{array}
\]
Every such normal form is realized by a term. The only constant also
in this family is the constant $I$.
\end{theorem}

We prove the upper bound, uniqueness, and synthesis separately. The
finite enumeration below independently checks the classification, but
does not replace these arguments.

\subsection{Upper bound and mandatory images}

A unary term is either constant or admits a live provenance
$(C,P)$ with nonempty $P\subseteq\operatorname{Sym}(I)$ such that
\[
 f(S)=C\cup\bigcup_{\pi\in P}\pi S\quad(|S|\geq2),\qquad
 f(\{i\})\subseteq C\cup\{\pi(i):\pi\in P\}.
\]
This follows by the same induction as Lemma~\ref{lem:provenance}, keeping
globally constant children separate. A constant owner singleton absorbs
globally; otherwise a live child remains live. In particular, a live
function preserves $I$ and maps ambiguous inputs to ambiguous outputs.

The induction also gives the following mandatory-image fact:
\begin{equation}\label{eq:mandatory}
 k\notin C\text{ and }\pi(i)=k\text{ for every }\pi\in P
       \quad\Longrightarrow\quad k\in f(\{i\}).
\end{equation}
At an owned node, every live child satisfying this hypothesis contains
$k$ at input $\{i\}$. If the owner differs from $k$, such a child cannot
be its singleton; a constant child cannot be that owner singleton
either, since the whole expression would have been globally constant.
Thus the node preserves $k$. If the owner is $k$, membership of $k$ is
disjunction and is preserved directly. Role switching transports this
argument; the variable supplies its base case.

For an output role $k$, let $A_k=\{\pi^{-1}(k):\pi\in P\}$. This set
is nonempty. On ambiguous inputs, membership of $k$ is either constantly
true, or the test $S\cap A_k\ne\varnothing$. With three roles, if
$|A_k|\geq2$, that test is always true. Hence every nonconstant output
coordinate is exactly one input membership bit, say $k\in f(S)$ iff
$i\in S$. Record $p(i)=k$. The recorded sources are distinct, because
every permutation in $P$ must send each such source to its recorded
target. Thus $p$ is a partial permutation, and the other coordinates
form $C_p$.

The atomic upper bound above excludes each variable output coordinate
from all rows except its source row. The mandatory-image
fact~\eqref{eq:mandatory} includes it in that row. These are precisely
the bounds in Theorem~\ref{thm:unary}. The partial permutation is
uniquely determined by its action on the three pairs: a variable output
coordinate is absent exactly on the pair omitting its source. Atomic
rows are actual values, so they too are unique.

\subsection{Synthesis of every normal form}

Use the already established reduct construction
\[
 R(C,E)(S)=\begin{cases}\{i\}\cup E_i,&S=\{i\},\\
 S\cup C,&|S|\geq2,\end{cases}\qquad E_i\subseteq C\setminus\{i\}.
\]
It uses only singleton constants and owned choice. Its generic syntactic
synthesis is now proved in \texttt{Synthesis.lean}.

\emph{Rank three.} The partial permutation is total, $C_p$ is empty,
and every row is its prescribed singleton. Use that role permutation.

\emph{Rank two.} Extend $p$ to a total permutation $\pi$. Its remaining
output is the single element of $C_p$. Apply a reduct function $R(C_p,E)$
after $\pi(x)$, choosing each row's optional extra winner as prescribed.
The remaining input row is the singleton in $C_p$. This realizes all
four row choices for each rank-two partial permutation.

\emph{Rank one.} By permuting the input, normalize $p$ to $p(a)=a$ and
$C_p=\{b,c\}$. First choose any desired nonempty rows $D_b,D_c$ within
$\{b,c\}$ while fixing the row $a$. On inputs $b,c$, owned choices by
$b,c$ are Boolean disjunction and conjunction, and the transposition
$(b\ c)$ is Boolean negation. Relative constants are
$x\ow b(b\ c)(x)$ and $x\ow c(b\ c)(x)$. Every expression so built
still sends $\{a\}$ to $\{a\}$. Taking unions of two Boolean functions
realizes any prescribed pair of nonempty rows $D_b,D_c$.

Precompose this expression with $R(C_p,0)$, which fixes all atomic inputs
and sends ambiguous inputs to $C_p$ or $I$. The Boolean expressions
above fix both those sets, so the resulting ambiguous behavior is
exactly $F_p$. Finally, to add any prescribed $Z\subseteq C_p$ to the
$a$ row, postcompose with $R(Z,E)$ having $E_a=Z$ and the other rows
empty. It preserves every nonempty subset of $C_p$, as well as $I$.
Thus all four choices for the $a$ row and all nine choices for the other
two rows are realized independently.

\emph{Rank zero.} The ambiguous output is always $I$ and the three
atomic rows may be any nonempty sets. The term
\[
 x\ow a(a\ b)(x)
\]
sends singleton $b$ to $a$ and fixes the other singleton roles. It
preserves $I$. These elementary mergers, together with permutations,
realize every map $h:I\to I$: merge each input into a chosen
representative of its $h$-fiber, then permute the distinct representatives
to their desired images. Representatives are never merged, so the
construction has the asserted effect. Union at most three such maps
to realize arbitrary nonempty atomic rows, and precompose with $R(I,0)$.
All ambiguous inputs then become $I$, while atomic rows are unchanged.

This proves synthesis. Larger ground constants give the constant
functions. Any nonconstant live function takes $I$ to $I$, so the only
constant overlap is $I$ itself.

\begin{corollary}[Exact count]
There are exactly $751$ unary term operations in the three-role full
signature. Of these, $745$ preserve $I$.
\end{corollary}
\begin{proof}
For rank $r$, the number of partial permutations is
$\binom3r^2r!$. The number of row choices for each is:
\[
\begin{array}{c|rrrr}
 r&0&1&2&3\\\hline
 \text{partial permutations}&1&9&18&6\\
 \text{row choices each}&343&36&4&1
\end{array}
\]
Hence there are $343+324+72+6=745$ normal forms of
type~\eqref{eq:unary}. Add the six constant functions other than $I$.
\end{proof}

\section{Context composition and the symmetric inverse monoid}

The partial permutations of $I$, including the empty map, form the
standard symmetric inverse monoid $\mathcal I_I$ under partial
composition~\cite{East2006}. For three roles its size is
$\sum_{r=0}^3\binom3r^2r!=34$.

\begin{theorem}[An ambiguity-action quotient]\label{thm:monoid}
The $745$ unary term operations preserving $I$ form a monoid under
composition. Sending the normal form $(p,D)$ to $p$ is a surjective
monoid homomorphism onto $\mathcal I_I$.
\end{theorem}
\begin{proof}
Every ambiguous input has an ambiguous output, and direct calculation
gives $F_p\circ F_q=F_{p\circ q}$. Indeed, an output coordinate either
is fixed present, or follows a unique input coordinate through each of
the two partial permutations. The atomic rows of a composite are
\[
 D^{f\circ g}_i=\begin{cases}
 D^f_j,&D^g_i=\{j\},\\
 F_{p_f}(D^g_i),&|D^g_i|\geq2.
 \end{cases}
\]
These formulas give the exact composite. Every $p$ has an admissible row
family and is realized by the synthesis, so the homomorphism is onto.
\end{proof}

\begin{corollary}[All reversible unary operations]
The invertible unary term operations are exactly the six role
permutations.
\end{corollary}
\begin{proof}
An invertible function fixes $I$: every other possibility is a constant.
Its partial-permutation image must be a unit, hence a total permutation.
Rank three leaves no choice of atomic rows. Conversely every role
permutation has its inverse as a term.
\end{proof}

Thus the familiar inverse monoid is a quotient governing ambiguity,
not the whole algebra of unary contexts. Its fibers retain the permitted
atomic rows. This is genuine context composition with internal
role-switching; it does not use uniform-profile selection, whose lift
failed to preserve substitution in the preceding note.

\section{Mechanization and finite checks}

\texttt{Synthesis.lean} now proves, generically for a finite list covering
all roles and three distinct owners:
\begin{itemize}
\item generation of every nonempty constant from singleton leaves;
\item synthesis of every reduct unary normal form;
\item the singleton separators and ambiguous-winner adders;
\item necessity of both conditions in the reduct's complete-test theorem;
\item correspondence between term instantiation and full-strategy SPE
winner sets, including realization of every nonempty winner set.
\end{itemize}
The proof imports only \texttt{Std} through the preceding modules, uses
Lean 4.34.1, and has no admitted facts or added axioms. The general
role-switching provenance theorem, sharp pair theorem, and
three-role unary classification remain mathematical proofs in this note.

\texttt{check\_symmetric.py} performs two independent constructions of
the three-role unary functions. It exhausts closure under all owned
operations and role permutations, and separately synthesizes actual
syntax for every proposed normal form. The resulting sets of $751$
functions agree exactly. It also checks all $34^2$ partial-permutation
composition cases. Singleton and pair separators are checked at
$K=3,4,5,6$, and the multivariable provenance representation is compared
with full evaluations of binary and ternary examples. These finite
checks support the written arguments; they are not generic Lean proofs.

\section{The finite-basis route: valid hypothesis, unresolved bound}

For each fixed $K$ the signature is finite. Every $\star_r$ is already
a semilattice operation, so the generated variety is congruence
meet-semidistributive. This follows directly from the semilattice
witnesses preceding Willard's Theorem 2.1\cite{Willard2000}; it does not
require introducing neutral union as a primitive.

Willard's Theorem 4.3 consequently applies \emph{if} the generated
variety has a finite residual bound. That means a uniform finite bound
on all its subdirectly irreducible members, including those arising
from subalgebras and homomorphic images of arbitrary powers. The
congruence chain of the generating algebra itself does not give that
bound. Neither its existence nor its failure is established here. The
optional script \texttt{probe\_residual.py} certifies selected
subdirectly irreducible quotients of a specified small subpower; such
finite experiments do not establish a residual bound.

This leaves two concrete research targets: determine the residual bound,
or extend the term-function characterization to more variables and more
roles. The three-role unary result is a complete classification in the
original role-switching signature; the optimal pair theorem is uniform
in $K$ and arity. Neither is yet a finite equational axiomatization.
The standard ingredients---Boolean functions, partial permutations,
semilattice witnesses, and uniform matroids---are reused directly.
Priority of these particular owned-choice results remains unestablished.

\bibliographystyle{plain}
\bibliography{references}
\end{document}
